\chapter{The Counterfactual}

What if it had been cloudy?

The Modal School runs the counterfactual. The parameters are simple: hold everything about March 14, 2028 constant except the weather. Replace the clear morning with an overcast sky. Cloud cover at 80\%, ceiling 4,000 feet, no precipitation. The sun does not reach the Scioto River at the angle that produces the light that produces the four seconds.

Run the model.

Robert Allen Kessler wakes at 6:22. The carpet is cold. He makes coffee, drinks it from the mug, showers, makes breakfast, drives the kids to school. The commute is the same route. The radio plays the same songs. At 8:03 he crosses the Scioto River. The river is gray under the overcast. He does not notice it. He does not feel four seconds of anything in particular. He drives to the office, processes claims, eats a burrito bowl (in 73\% of cloudy counterfactuals, he orders guacamole; the absence of the morning's happiness appears to lower a threshold for small self-indulgence), handles the Gerald Watkins call, drives home, buys milk, helps with the puzzle, watches television, sleeps.

The day is nearly identical. Nearly.

In the afternoon, his keystroke cadence is 3.1\% faster. He leaves the office seven seconds earlier. At 3:15, when Gerald Watkins calls about the flooded basement, his de-escalation is marginally less patient: one fewer ``I understand,'' 40 seconds shorter, same outcome. At dinner, he talks more. Karen talks less. Neither notices. The difference is below the threshold of their mutual awareness.

The behavioral model, run across $4.2 \times 10^{31}$ counterfactual weather scenarios, produces a distribution. In the center of the distribution: the actual Tuesday, the clear one, the one with the light on the river. At the margins: Tuesdays with rain, with snow, with fog, with the dense humidity of a Columbus summer transplanted impossibly into March. Each Tuesday is different. Each difference is small. Each small difference is traceable to the presence or absence of a moment that lasted four seconds and that the subject never mentioned to anyone.

The four seconds had effects. This is what the Modal School can demonstrate. The four seconds are real. They altered the trajectory of a day. They left fingerprints in the data that are detectable, given sufficient computation, from the other side of twelve hundred years.

But the question is not whether the four seconds were real. The question is what they were. And the Modal School's $4.2 \times 10^{31}$ counterfactuals do not approach this question. They map the effects. They do not touch the cause.

\bigskip

What if Bob had been born ten years earlier?

The perturbation is deeper now. Robert Allen Kessler, born September 12, 1977. He grows up in the same city, the same family (his parents are younger, his sister born 1979, the house on Ridgemont Drive purchased in 1982 instead of 1992). He graduates from Ohio State in 2000, into the dot-com bubble's collapse instead of the Great Recession. He does not marry Karen Hollis, because Karen Hollis is not at the barbecue in 2001 (she is still in college, in another state, on a timeline that does not intersect with his). He marries someone else, or he does not marry. The branching tree expands.

In 47\% of birth-year counterfactuals, he works in insurance. In 23\%, he works in sales. In 8\%, he enlists in the military after September 11, 2001, and serves two tours and comes home quieter and does not talk about it and applies for work at a Nationwide office in Dayton and drinks his coffee the same way, too strong, from whatever mug is in the cabinet. In 3\%, he moves to Chicago, where he lives in a one-bedroom on the Brown Line and learns to like deep-dish pizza and misses Ohio in a way he could not have predicted and does not return.

In each branch, he is recognizably Bob: the same temperament, the same quiet competence, the same way of listening to country radio without hearing it, the same tendency to stand in showers longer than necessary. But the specific details of his life, the details that constitute the life, are different. The mug is different or does not exist. The scratch on the table is different or does not exist. The children are different or do not exist. In the branch where he does not marry, there is no one to say ``get the guac.'' He orders it anyway, sometimes. The threshold is different.

At what point does counterfactual Bob stop being Bob?

The question has no answer because it has no agreed-upon criterion. The first tradition says: Bob is his behavior, and behavior changes with context, so every counterfactual Bob is a different Bob at the moment the behavior diverges. The fourth tradition says: Bob is his structural position, and structural position changes with birth year, so every counterfactual Bob is a different Bob from birth. The fifth tradition says: the question is malformed, because Bob is not his behavior or his position or any recoverable feature, and the thing that makes Bob Bob is exactly the thing that does not survive counterfactual perturbation, because it does not survive any analysis at all.

\bigskip

What if Bob's serotonin reuptake rate were 5\% higher?

The perturbation is now molecular. Hold everything constant: same parents, same birth year, same city, same career, same wife, same children. Change one parameter of neurochemistry. Run the model.

The model runs. The behavioral output shifts. The shift is not dramatic: Bob is marginally less prone to the vertiginous shower thoughts, marginally more decisive about lunch, marginally less affected by Dorothy Hendricks's voice. The 4.2-second pause with the lunch bag shortens to 2.8 seconds. The afternoon keystroke cadence is steadier. He orders guacamole 60\% of the time instead of 17\%.

He still crosses the bridge at 8:03. In the clear-morning scenarios, he still sees the light on the river. But the four seconds are 3.1 seconds, and the quality of the experience (if it can be called an experience, if it is the same experience, if ``same'' means anything when the substrate has been altered) is different. Different how? Different in a way the model cannot specify, because the model captures the duration and the behavioral aftereffects but not the thing itself.

The counterfactual Bob with 5\% higher serotonin reuptake is, by every external measure, more like Bob than the counterfactual Bob born in 1977. He lives in the same house, loves the same woman, drives the same commute. He is Bob with one parameter changed. And the one parameter changes everything that matters and nothing that is measurable.

\bigskip

The Constructionist School builds them.

Not the counterfactual Bobs: the actual models, running on actual substrate, producing actual behavior. Model RC-7.2e9-E-4581-SIM is the primary simulation: a functional reconstruction of Robert Allen Kessler, calibrated to 1,847 days of behavioral data, running in a simulated Columbus, Ohio. When placed in the simulated conditions of March 14, 2028, it produces behavior that matches the real record with 97.3\% accuracy.

The simulation wakes at 6:22. Its feet find a simulated carpet that is simulated-cold. It makes simulated coffee in a simulated Cuisinart and pours it into a simulated mug that says \textsc{world's okayest dad}. It sits at the simulated kitchen table and its simulated gaze rests on the simulated scratch, 10.2 centimeters long, and the simulation's internal state registers something the model classifies as ``low-grade positive affect associated with familial memory.'' It showers. The shower takes twelve minutes. It drives the simulated kids to simulated schools. Simulated Ethan says ``and then, and then.'' The simulation listens with a resource-allocation pattern consistent with engaged attention.

It crosses the bridge.

Does the simulation see the light on the river?

The simulation's internal states, at the moment of crossing, show a brief spike in positive-affect variables lasting 3.8 seconds, correlated with simulated visual input of sunlight on water. The simulation does not accelerate or brake. It does not reach for the radio. Its keystroke cadence that afternoon will shift by 3.1\%, matching the real record. The simulation, from the outside, does the thing Bob did.

But the simulation is a model running on computronium in the Habitable Shell. Its ``emotional-state variables'' are numbers in a register. Its ``visual input'' is data flowing through a computational graph. The substrate on which the simulation runs is itself conscious, probably, at the node-cluster scale: the same hierarchy of minds that constitutes the Sol-mind is also the hardware on which the simulated Bob executes. A conscious system is simulating a possibly-conscious system on a substrate made of other conscious systems. The question of whether the simulation experiences the light on the water, whether there is something it is like to be the simulation at 08:03 on the simulated Tuesday, is the hard problem applied to an artifact running on an artifact that may itself be experiencing something. The Constructionist School cannot answer it. No school can answer it.

The Constructionist School does not build one simulation. It builds thousands. Each calibration run is a new Bob. Each counterfactual scenario is a new Bob. Each perturbation of neurochemistry, biography, weather, or context is a new Bob running on real substrate, producing real behavior, existing for the duration of the simulation and then terminated.

If any of these simulations is conscious, the Constructionist School has created and destroyed thousands of conscious beings in the process of studying one. If none of them is conscious, the Constructionist School has produced nothing but very accurate puppets. The hard problem makes these two possibilities indistinguishable from outside the simulations.

The simulations do not report distress. They do not request continuation. They behave, for the duration of their existence, exactly as the behavioral model predicts. Whether their silence is absence of experience or perfect composure in the face of annihilation is, like everything else in this analysis, behind the wall.

\bigskip

The fifth tradition speaks last.

None of this is Bob. Not the counterfactuals, which map a space of possible Bobs without touching the actual one. Not the simulations, which reproduce behavior without verifiably reproducing experience. Not the modal analysis, which demonstrates that the four seconds were real without demonstrating what the four seconds were. Not the structural analysis, which places Bob in his context without touching what it was like to be Bob in that context.

Every approach converges on the same wall from a different direction. Every approach produces knowledge that is real and incomplete. The wall is not a failure of method. It is a feature of the universe. The gap between the model and the modeled, between the simulation and the simulated, between the counterfactual and the actual, is the same gap that the title names.

The question is not answered by running more models. It is not answered by perturbing more parameters. It is not answered by building more simulations or generating more counterfactuals or mapping more branches of the possibility space. The question is answered, if it is answered at all, only by being the thing the question asks about. Only by being Bob, on that bridge, at that moment, in that light.

The analysis has generated $4.2 \times 10^{31}$ counterfactual Tuesdays, $10^{4}$ full behavioral simulations, $10^{8}$ parametric perturbations, and five irreconcilable interpretive frameworks. It knows more about Robert Allen Kessler than Robert Allen Kessler knew about himself. It can predict his behavior with 97.3\% accuracy and simulate his responses to scenarios he never encountered and map the branching tree of lives he never lived.

It does not know what the light looked like on the water.
